% Generated draft: requires mathematical and source review.
\documentclass[11pt]{article}
\usepackage[margin=25mm]{geometry}
\usepackage{fontspec}
\setmainfont{Latin Modern Roman}
\usepackage{amsmath,amssymb,mathtools,unicode-math}
\setmathfont{Latin Modern Math}
\usepackage{enumitem,needspace,xurl,hyperref}
\setlength{\parindent}{0pt}
\setlength{\parskip}{5pt}
\setlength{\emergencystretch}{3em}
\newcommand{\N}{\mathbb{N}}
\newcommand{\Z}{\mathbb{Z}}
\newcommand{\Q}{\mathbb{Q}}
\newcommand{\R}{\mathbb{R}}
\newcommand{\C}{\mathbb{C}}
\newcommand{\F}{\mathbb{F}}
\newcommand{\A}{\mathbb{A}}
\newcommand{\PP}{\mathbb{P}}
\newcommand{\E}{\mathbb{E}}
\newcommand{\eps}{\varepsilon}
\newcommand{\dd}{\,\mathrm{d}}
\begin{document}
\section*{UnsolvedMath ID: 700890}
\textbf{Source problem code:} catalogue entry
\subsection{Definitions and mathematical statement}
A formation of finite groups is a nonempty class closed under isomorphisms, homomorphic images, and finite subdirect products. It is hereditary if it is closed under taking subgroups, and soluble if all its members are soluble. Here local is understood in the usual chief-factor sense: there are local values \(f(p)\), each a formation or the empty class, indexed by primes, such that a finite group \(H\) belongs to \(\mathfrak F\) precisely when
\[
 H/C_H(A/B)\in f(p)
\]
for every chief factor \(A/B\) of \(H\) and every prime \(p\) dividing \(|A/B|\). Here
\[
 C_H(A/B)=\{h\in H:[a,h]\in B\text{ for every }a\in A\}.
\]
A chief factor is a section \(A/B\), with \(B<A\) normal in \(H\), having no normal subgroup of \(H\) strictly between \(B\) and \(A\).

Let \(\mathfrak N\) be the formation of all finite nilpotent groups. For a finite soluble group \(G\), its nilpotent length \(l(G)\) is the least integer \(r\geq 0\) for which there is a series
\[
 1=G_0\trianglelefteq G_1\trianglelefteq\cdots\trianglelefteq G_r=G
\]
with each \(G_i\) normal in \(G\) and each factor \(G_i/G_{i-1}\) nilpotent. In particular, \(l(1)=0\). For \(k\geq 0\), interpret
\[
 \mathfrak N^k=\{G:G\text{ is finite soluble and }l(G)\leq k\},
 \qquad \mathfrak N^0=\{1\}.
\]
For a class \(\mathfrak X\), a finite group \(G\) is minimal non-\(\mathfrak X\) if \(G\notin\mathfrak X\) and every proper subgroup of \(G\) belongs to \(\mathfrak X\).

The supplied record does not define ``radical.'' In the statement below it is interpreted in the conventional Fitting-class sense: the class is closed under normal subgroups and under products of normal members in a common finite group. For a hereditary formation, the normal-subgroup condition already holds. Thus the remaining condition is
\[
 \forall H\text{ finite},\quad
 A,B\trianglelefteq H,\ A,B\in\mathfrak F
 \ \Longrightarrow\ AB\in\mathfrak F.
\]
This interpretation of the source terminology requires confirmation in mathematical review.

\textbf{Requested statement (under these conventions).}
Let \(\mathfrak F\) be a soluble local hereditary formation of finite groups. Suppose that for every finite soluble group \(G\),
\[
 \left(G\notin\mathfrak F\ \text{ and }\ \forall K<G,\ K\in\mathfrak F\right)
 \Longrightarrow
 \left(G\notin\mathfrak N^{\,l(G)-1}\ \text{ and }\
 \forall K<G,\ K\in\mathfrak N^{\,l(G)-1}\right).
\]
Then \(\mathfrak F\) is radical in the sense specified above. Every group \(G\) covered by the antecedent is nontrivial, so \(l(G)\geq 1\) and the exponent \(l(G)-1\) is defined. Equivalently, the hypothesis requires every proper subgroup \(K<G\) of every finite soluble minimal non-\(\mathfrak F\) group to satisfy \(l(K)\leq l(G)-1\).
\subsection{Short English statement}
For a soluble local hereditary formation \(\mathfrak F\), assume that every finite soluble group minimal outside \(\mathfrak F\) has all proper subgroups of strictly smaller nilpotent length. The requested conclusion is that \(\mathfrak F\) is radical, interpreted here as closure under products of normal members in a common finite group. The source terminology needs confirmation. The supplied registry labels this problem open; that label is not an independently verified assessment.
\subsection{Sources}
\begin{itemize}
\item Supplied local UnsolvedMath registry record: permanent ID \(700890\), exact source problem code \texttt{catalogue entry}, title ``Kourovka Notebook Problem 12.72.'' The record supplies the status ``open.''
\item \emph{The Kourovka Notebook (Unsolved Problems in Group Theory)}, original problem 12.72, proposed by V. N. Semenchuk in the 12th issue (1992), as attributed in the supplied record.
\item The supplied background identifies the 21st issue edition (2026), \texttt{arXiv:1401.0300v48}. The supplied collection description attributes editing of Issues 1--21 (1965--2026) to E. I. Khukhro and V. D. Mazurov. These bibliographic details are supplied attribution; the referenced edition has not been independently inspected.
\end{itemize}
\end{document}
